\documentclass[12pt]{article}
\usepackage[utf8]{inputenc}
\usepackage[T1]{fontenc}
\usepackage[margin=1in]{geometry}
\usepackage{setspace}
\usepackage{xcolor}
\usepackage[colorlinks=true, urlcolor=blue!60!black, linkcolor=blue!60!black]{hyperref}
\setlength{\parindent}{0em}
\setlength{\parskip}{0.6em}
\renewcommand{\baselinestretch}{1.0}
\begin{document}
\begin{center}
{\LARGE \textit{Hernandez v.\ Texas}, 347 U.S.\ 475 (1954)\par}
\vspace{0.5em}
{\large Case Brief --- Two-Page Edition\par}
\end{center}

\section{Memory Jogger}
Equal protection is not limited to a Black--white ``two-class theory'': Texas could not systematically exclude Mexican Americans from juries (\href{https://supreme.justia.com/cases/federal/us/347/475/}{Justia}).

\section{Detailed Case Facts}
Pete Hernandez, a Mexican-American farmworker, was indicted and convicted of murder by all-Anglo juries in Jackson County, Texas. Texas stipulated that no person with a Mexican or Latin American name had served on any jury there in 25 years, out of more than 6,000 jurors, though qualified Mexican Americans lived in the county. Local segregation was pervasive, down to a courthouse restroom marked ``Colored Men'' and ``Hombres Aqui.''

\section{Procedural History}
The Texas Court of Criminal Appeals affirmed, reasoning that the Fourteenth Amendment recognized only white and Negro classes and that Mexican Americans were white.

\section{Judicial Votes}
Reversed, 9--0. Chief Justice Warren wrote for a unanimous Court.

\section{Holding}
The Equal Protection Clause bars excluding eligible jurors because of ancestry or national origin. A challenger must prove that the group is a distinct class in the local community and make a prima facie case of systematic exclusion under \textit{Norris v. Alabama}; general denials of intent do not rebut it.

\section{Analysis of Opinions}
\textbf{Majority:} Groups needing protection change as community prejudices change. Testimony, segregated schooling, a ``No Mexicans Served'' sign, and the courthouse restroom proved that Mexican Americans were a class apart, and Spanish surnames identified them as readily as skin color identifies race. With qualified Mexican Americans in the county and none chosen in 25 years, chance could not explain the result, and the commissioners' professions of good faith could not overcome the proof. No Justice wrote separately.

\section{Examples: Future Applications}
\textbf{Same-Side:} a county that has called no Hmong-American jurors in 20 years despite a large, segregated Hmong community violates equal protection. \textbf{Opposite-Side:} a defendant whose randomly drawn jury happens to include no Mexican Americans has no claim, because there is no right to proportional representation. \textbf{Fence-Sitter:} Mexican Americans are 30\% of residents but 4\% of grand jurors---underrepresentation without total exclusion.

\section{Critique}
Ian Haney L\'{o}pez called it ``the single most insightful Supreme Court opinion on race ever handed down,'' yet scholars note that its county-by-county proof of class status burdened Latino litigants and left the racial binary intact. Formalists find its willingness to infer discrimination without conscious intent hard to square with \textit{Washington v. Davis}.

\section{Key Quotations}
\begin{quote}
``Throughout our history, differences in race and color have defined easily identifiable groups which have at times required the aid of the courts in securing equal treatment under the laws. But community prejudices are not static, and, from time to time, other differences from the community norm may define other groups which need the same protection.''
\end{quote}
\noindent --- Chief Justice Warren, for the Court, 347 U.S. at 478.

\end{document}
